% SEGMENT OLP-0658-S01
% Part: proof-theory
% Chapter: cut-elimination
% Section: interpolation

\documentclass[../../../include/open-logic-section]{subfiles}

\begin{document}

\olfileid{pt}{cut}{itp}

\olsection{மாஎஹாராவின் துணைத்தேற்றமும் கிரெய்கின் இடையீட்டுத் தேற்றமும்}

வில்லியம் கிரெய்கின் இடையீட்டுத் தேற்றத்தின்படி, $!A \Entails !B$
எனில், $!A \Entails !C$ மற்றும் $!C \Entails !B$ ஆகிய இரண்டும்
நிறைவேறுமாறு !!a{sentence}~$!C$ ஒன்று உள்ளது. இதுவே
\emph{இடையீட்டி}. முக்கியமாக, $!C$ இல் வரும் !!{constant}s மற்றும்
!!{predicate}s அனைத்தும் $!A$ மற்றும் $!B$ இரண்டிலும் வருமாறு
$!C$ ஐத் தேர்ந்தெடுக்கலாம். இங்கு !!{function}s இல்லை எனக்
கொள்கிறோம். இந்தப் பொதுக் குறியீட்டு நிபந்தனை இல்லாவிட்டால்,
$!A$ அல்லது $!B$ ஐயே எளிதாக இடையீட்டியாகக் கொள்ளலாம்.

இத்தேற்றம் செவ்வியல் முதல்நிலைத் தருக்கத்திற்குப் பொருந்துகிறது;
கண்டறியப்பட்ட முதல்நிலைத் தருக்கத்தின் ``கடைசி முக்கியமான பண்பு''
என்றும் இது அழைக்கப்பட்டுள்ளது. மாதிரிக் கோட்பாட்டிலும் தானியங்கித்
தருக்கவாதத்திலும் இதற்குப் பயன்கள் உள்ளன. அறிவியல் தத்துவத்தில் ஒரு
கோட்பாட்டை எந்த அடிப்படைக் குறியீடுகளால் அடிக்கோளாக்கலாம் அல்லது
முடியாது என்பதை ஆய்வு செய்வதே இதன் தொடக்க நோக்கமும் பயன்பாடுமாக
இருந்தது. ஒரு கோட்பாடு ஒரு கருத்தை மறைமுகமாக வரையறுக்க இயன்றால்,
அதை வெளிப்படையாகவும் வரையறுக்க இயலும் எனும் பெத்தின்
வரையறுக்கத்தக்கமைத் தேற்றமும் இதிலிருந்து பின்தொடர்கிறது.

கணிதத் தருக்கத்தின் பல முடிவுகளைப் போலவே இடையீட்டுத் தேற்றத்தையும்
மாதிரிக் கோட்பாட்டு முறைகளாலும் நிறுவல்-கோட்பாட்டு முறைகளாலும்
நிறுவலாம். நிறுவல்-கோட்பாட்டு முறையின் சிறப்பு, இடையீட்டி இருக்கிறது
எனக் காட்டுவதோடு நிற்காமல், $!A \Entails !B$ என்பதற்கான
!!a{proof} இலிருந்து அதை கணக்கிடும் படிமுறையையும் தருவதாகும்;
எடுத்துக்காட்டாக, $!A \Sequent !B$ என்பதற்கான~\Log{G3c}
!!a{proof} இலிருந்து கணக்கிடலாம்.

$\Lang L(!A)$ என்பது $!A$ இல் வரும் !!{constant}s மற்றும்
!!{predicate}s இன் கணம் என்க; மேலும்
$\Lang L(\Gamma) = \bigcup \Setabs{\Lang L(!A)}{!A \in
\Gamma}$ என்க. இப்பிரிவு முழுவதும் !!{function}s இல்லாத
!!{formula}s ஐயே கருதுகிறோம்.

% SEGMENT OLP-0658-S02
\begin{thm}[கிரெய்கின் இடையீட்டுத் தேற்றம்]\ollabel{thm:interpolation}
$!A$ இன் !!{language}~$\Lang L(!A)$ என்றும் $!B$ இன்
மொழி~$\Lang L(!B)$ என்றும் கொள்க. $\Log{G3c} \Proves
!A \Sequent !B$ எனில், $\Lang L(!C) \subseteq \Lang L(!A)
\cap \Lang L(!B)$ ஆகும் !!{language} இல் !!a{formula}~$!C$
ஒன்று உள்ளது; மேலும் $\Log{G3c} \Proves !A \Sequent !C$
மற்றும் $\Log{G3c} \Proves !C \Sequent !B$.
\end{thm}

~\Log{G3c} இன் !!{proof}s இன் அமைப்பைப் பயன்படுத்துவதற்காக,
இடையீட்டுத் தேற்றத்தைப் பொதுமைப்படுத்தி நிறுவுகிறோம். ஒரு தொடரணியின்
முன்பக்கத்தையும் பின்பக்கத்தையும் தலா இரு பகுதிகளாகப் பிரித்துக்
கொள்க. $\maeh{\Gamma_1}{\Gamma_2} \Sequent
\maeh{\Delta_1}{\Delta_2}$ என்ற குறியில், $\Gamma_1$ உம்
$\Delta_1$ உம் முதல் பகுதிக்கும், $\Gamma_2$ உம் $\Delta_2$ உம்
இரண்டாம் பகுதிக்கும் உரியவை. பின்வரும் துணைத்தேற்றத்திலிருந்து
இடையீட்டுத் தேற்றம் உடனே பெறப்படும்.

\begin{lem}[மாஎஹாராவின் துணைத்தேற்றம்]\ollabel{lem:maehara}
$\Log{G3c} \Proves \maeh{\Gamma_1}{\Gamma_2} \Sequent
\maeh{\Delta_1}{\Delta_2}$ எனில்,
$\Lang L(!C) \subseteq \Lang L(\Gamma_1,
\Delta_1) \cap \Lang L(\Gamma_2, \Delta_2)$ ஆகும்
!!{language} இல் !!a{formula}~$!C$ ஒன்று உள்ளது; மேலும்
$\Log{G3c} \Proves \Gamma_1 \Sequent \Delta_1, !C$
மற்றும் $\Log{G3c} \Proves !C, \Gamma_2 \Sequent \Delta_2$.
\end{lem}

% SEGMENT OLP-0658-S03
\begin{proof}
  $\pi \Proves \maeh{\Gamma_1}{\Gamma_2} \Sequent
  \maeh{\Delta_1}{\Delta_2}$ என்று கொள்க. $n=\pheight{\pi}$
  மீது தொகுத்தறிந்து கூற்றை நிறுவுகிறோம்.

  $n=0$ எனில், $\maeh{\Gamma_1}{\Gamma_2} \Sequent
  \maeh{\Delta_1}{\Delta_2}$ ஓர் அடிக்கோள். அப்போது ஓர் அணு
  !!{formula}~$!A$, $\Gamma_1, \Gamma_2$ இலும்
  $\Delta_1, \Delta_2$ இலும் வருகிறது. இடப்புறத்தில் $!A$
  எந்தப் பகுதியில் வருகிறது என்பதையும் வலப்புறத்தில் எந்தப் பகுதியில்
  வருகிறது என்பதையும் பொறுத்து நான்கு நேர்வுகள் உள்ளன.
  \begin{enumerate}
    \item $!A \in \Gamma_1$ மற்றும் $!A \in \Delta_1$.
    $\Gamma_1 \Sequent \Delta_1, !C$ மற்றும்
    $!C, \Gamma_2 \Sequent \Delta_2$ ஆகிய இரண்டும்
    !!{provable} ஆகும் $!C$ தேவை. $!A$ இடப்புறமும் வலப்புறமும்
    வருவதால், $!C$ எதுவாயினும் முதலாவது தொடரணி ஏற்கெனவே
    அடிக்கோள். $!C, \Gamma_2 \Sequent \Delta_2$ என்ற
    இரண்டாவது தொடரணி !!{provable} என்பதை உறுதிப்படுத்துவதற்கு
    $!C \ident \lfalse$ எனத் தேர்ந்தெடுக்கிறோம்.
    \item $!A \in \Gamma_1$ மற்றும் $!A \in \Delta_2$.
    அப்போது $\Gamma_1 \Sequent \Delta_1, !A$ மற்றும்
    $!A, \Gamma_2 \Sequent \Delta_2$ இரண்டும் அடிக்கோள்கள்;
    எனவே $!C \ident !A$ எனலாம்.
    \item $!A \in \Gamma_2$ மற்றும் $!A \in \Delta_1$.
    $!A, \Gamma_1 \Sequent \Delta_1$ மற்றும்
    $\Gamma_2 \Sequent \Delta_2, !A$ அடிக்கோள்களாக
    இருக்கும்; ஆனால் இடையீட்டிக்குத் தேவையான பக்கங்களுக்கு
    நேர்மாறாக $!A$ அமைந்துள்ளது. ஆகவே $!A$ இடையீட்டி அல்ல.
    இருப்பினும் $\RightR{\lnot}$ மற்றும் $\LeftR{\lnot}$
    விதிகளால் $\Gamma_1 \Sequent \Delta_1, \lnot!A$
    மற்றும் $\lnot !A, \Gamma_2 \Sequent \Delta_2$
    ஆகியவற்றை !!{prove} செய்யலாம்.
    \item $!A \in \Gamma_2$ மற்றும் $!A \in \Delta_2$.
    பயிற்சி.
  \end{enumerate}
  $\lfalse \in \Gamma_1$ அல்லது $\lfalse \in \Gamma_2$
  என்பதால் $\maeh{\Gamma_1}{\Gamma_2} \Sequent
  \maeh{\Delta_1}{\Delta_2}$ அடிக்கோளாகும் நேர்வுகளும்
  பயிற்சிகளாக விடப்படுகின்றன.

  $n>0$ எனில், $\pi$ தருக்க அனுமானத்தில் முடிகிறது. பயன்படுத்தப்பட்ட
  விதி எது, அதன் முதன்மை வாய்பாடு எந்தப் பகுதியைச் சேர்ந்தது என்று
  பிரித்து நேர்வுகளைப் பார்க்க வேண்டும். $\pheight{\pi_1}<n$
  ஆகும் ஒவ்வொரு !!{proof}~$\pi_1$ க்கும், $\pi_1$ இன் இறுதித் தொடரணி
  எவ்வாறு இரண்டாகப் பிரிக்கப்பட்டிருந்தாலும், துணைத்தேற்றம் பொருந்தும்
  என்பதே தொகுத்தறிதல் கருதுகோள். சில நேர்வுகளை விரிவாகப் பார்ப்போம்.
  \begin{enumerate}
% SEGMENT OLP-0658-S04
    \item $\pi$, முதன்மை வாய்பாடு~$\lnot !A$ உடைய
    $\RightR{\lnot}$ விதியில் முடிகிறது. $\lnot !A$ முதல் பகுதியிலா,
    இரண்டாம் பகுதியிலா உள்ளது என்பதைப் பொறுத்து $\pi$ பின்வரும்
    இரண்டில் ஒன்றில் முடிகிறது:
    \[
    \AxiomC{}
    \RightLabel{$\pi_1$}
    \Deduce$\maeh{!A, \Gamma_1}{\Gamma_2} \fCenter
      \maeh{\Delta_1}{\Delta_2}$
    \RightLabel{\RightR{\lnot}}
    \UnaryInf$\maeh{\Gamma_1}{\Gamma_2} \fCenter
      \maeh{\Delta_1, \lnot !A}{\Delta_2}$
    \DisplayProof
    \]
    அல்லது
    \[
    \AxiomC{}
    \RightLabel{$\pi_1$}
    \Deduce$\maeh{\Gamma_1}{!A, \Gamma_2} \fCenter
      \maeh{\Delta_1}{\Delta_2}$
    \RightLabel{\RightR{\lnot}}
    \UnaryInf$\maeh{\Gamma_1}{\Gamma_2} \fCenter
      \maeh{\Delta_1}{\Delta_2, \lnot !A}$
    \DisplayProof
    \]
    முதன்மை !!{formula}~$\lnot !A$ முதல் பகுதியில் இருந்தால்,
    முன்கோளில் அதன் செயல்படு !!{formula}~$!A$ ஐயும் அதே பகுதியில்
    வைத்துள்ளோம். இது நாம் செய்யக்கூடிய தேர்வு. செயல்படு வாய்பாட்டை
    மறுபகுதியில் வைத்தாலும் தொகுத்தறிதல் கருதுகோள் பொருந்தும்;
    ஆனால் அப்போது இடையீட்டியைக் கண்டுபிடிக்க முடியாது
    (பயிற்சி: முயன்று பார்க்க).

    முதலில் $\lnot !A$ முதல் பகுதியில் உள்ள நேர்வைக் கருதுக.
    தொகுத்தறிதல் கருதுகோளால் $\pi_1$ இன் இறுதித் தொடரணிக்கு
    இடையீட்டியான !!a{formula}~$!C_1$ ஒன்று உள்ளது; அதாவது,
    \begin{align*}
    !A, \Gamma_1 & \Sequent \Delta_1, !C_1 &&\text{மற்றும்} &
    !C_1, \Gamma_2 & \Sequent \Delta_2
    \intertext{ஆகிய இரண்டும் நிறுவத்தக்கவை. நமக்குத் தேவைப்படும்
    !!a{formula}~$!C$ பின்வரும் இரண்டையும் !!{provable} ஆக்க வேண்டும்:}
    \Gamma_1 & \Sequent \Delta_1, \lnot !A, !C &&\text{மற்றும்} &
    !C, \Gamma_2 & \Sequent \Delta_2
    \end{align*}
    $!C \ident !C_1$ எனவே கொள்ளலாம். வலப்புறத் தொடரணி
    !!{provable} என்று தொகுத்தறிதல் கருதுகோள் ஏற்கெனவே சொல்கிறது;
    இடப்புறத் தொடரணி, அதே கருதுகோளில் கிடைத்த தொடரணிக்கு
    $\RightR{\lnot}$ பயன்படுத்துவதால் கிடைக்கிறது. மேலும்
    $\Lang L(!C_1) \subseteq \Lang L(!A, \Gamma_1,
    \Delta_1) \cap \Lang L(\Gamma_2, \Delta_2)$.
    $\Lang L(\lnot !A) = \Lang L(!A)$ மற்றும்
    $\Lang L(!C) = \Lang L(!C_1)$ என்பதால்,
    $\Lang L(!C_1) \subseteq \Lang L(\Gamma_1, \Delta_1,
    \lnot !A) \cap \Lang L(\Gamma_2, \Delta_2)$ என்பதும்
    நிறைவேறுகிறது.

    இரண்டாம் நேர்வில் $\lnot !A$ இரண்டாம் பகுதியில் உள்ளது.
    முன்கோளின் செயல்படு !!{formula} ஐயும் இரண்டாம் பகுதியில் வைத்து,
    தொகுத்தறிதல் கருதுகோளைப் பயன்படுத்தினால்,
    \begin{align*}
    \Gamma_1 & \Sequent \Delta_1, !C_1 &&\text{மற்றும்} &
    !C_1, !A, \Gamma_2 & \Sequent \Delta_2
    \intertext{ஆகியவை நிறுவத்தக்கவை. இவற்றில் இரண்டாவதற்கு
    \RightR{\lnot} விதியைப் பயன்படுத்தினால் பின்வரும்
    தொடரணிகளுக்கு !!{proof}s கிடைக்கின்றன:}
    \Gamma_1 & \Sequent \Delta_1, !C_1 &&\text{மற்றும்} &
    !C_1, \Gamma_2 & \Sequent \Delta_2, \lnot !A
    \end{align*}
    $!C_1$ இடையீட்டி ஆக வேண்டியபோது !!{provable} ஆக வேண்டிய
    தொடரணிகள் இவையே; மீண்டும் $!C \ident !C_1$ என்கிறோம்.
    முன்போலவே $\Lang L(!C_1) \subseteq \Lang L(\Gamma_1,
    \Delta_1) \cap \Lang L(\Gamma_2, \Delta_2, \lnot !A)$.

% SEGMENT OLP-0658-S05
    \item !!{proof}~$\pi$, முதன்மை !!{formula}~$!A \lif !B$ உடைய
    $\RightR{\lif}$ விதியில் முடிகிறது. $!A \lif !B$ முதல்
    அல்லது இரண்டாம் பகுதியில் இருப்பதைப் பொறுத்து, $\pi$ பின்வரும்
    இரண்டில் ஒன்றில் முடிகிறது:
    \[
    \AxiomC{}
    \RightLabel{$\pi_1$}
    \Deduce$\maeh{!A, \Gamma_1}{\Gamma_2} \fCenter
      \maeh{\Delta_1, !B}{\Delta_2}$
    \RightLabel{\RightR{\lif}}
    \UnaryInf$\maeh{\Gamma_1}{\Gamma_2} \fCenter
      \maeh{\Delta_1, !A \lif !B}{\Delta_2}$
    \DisplayProof
    \]
    அல்லது
    \[
    \AxiomC{}
    \RightLabel{$\pi_1$}
    \Deduce$\maeh{\Gamma_1}{!A, \Gamma_2} \fCenter
      \maeh{\Delta_1}{\Delta_2, !B}$
    \RightLabel{\RightR{\lif}}
    \UnaryInf$\maeh{\Gamma_1}{\Gamma_2} \fCenter
      \maeh{\Delta_1}{\Delta_2, !A \lif !B}$
    \DisplayProof
    \]
    இங்கும் முடிவின் முதன்மை !!{formula} எந்தப் பகுதியில் உள்ளதோ,
    முன்கோளின் செயல்படு !!{formula}s ஐயும் அதே பகுதியில் வைத்துள்ளோம்.
    முதல் நேர்வில் தொகுத்தறிதல் கருதுகோள்,
    \begin{align*}
    !A, \Gamma_1 & \Sequent \Delta_1, !B, !C_1 &&\text{மற்றும்} & !C_1, \Gamma_2 & \Sequent \Delta_2
    \intertext{ஆகியவற்றுக்கு !!{proof}s தருகிறது. இடப்புறத் தொடரணி
    \RightR{\lif} விதிக்குப் பொருத்தமான முன்கோள். ஆகவே பின்வரும்
    தொடரணிகள் !!{provable}:}
    \Gamma_1 & \Sequent \Delta_1, !A \lif !B, !C_1 &&\text{மற்றும்} & !C_1, \Gamma_2 & \Sequent \Delta_2
    \end{align*}
    இதனால் $!C_1$, $\pi$ இன் இறுதித் தொடரணிக்கும் இடையீட்டி;
    மீண்டும் $!C \ident !C_1$ எனலாம். மற்ற நேர்வும் இதேபோல்
    கையாளப்படுகிறது.

% SEGMENT OLP-0658-S06
    \item !!{proof}~$\pi$, முதன்மை !!{formula}~$!A \lif !B$ உடைய
    $\LeftR{\lif}$ விதியில் முடிகிறது. $!A \lif !B$ முதல்
    பகுதியில் இருந்தால், $\pi$ பின்வருமாறு முடிகிறது:
    \[
    \AxiomC{}
    \RightLabel{$\pi_1$}
    \Deduce$\maeh{\Gamma_1}{\Gamma_2} \fCenter
      \maeh{\Delta_1, !A}{\Delta_2}$
    \AxiomC{}
    \RightLabel{$\pi_2$}
    \Deduce$\maeh{!B, \Gamma_1}{\Gamma_2} \fCenter
      \maeh{\Delta_1}{\Delta_2}$
    \BinaryInf$\maeh{!A \lif !B, \Gamma_1}{\Gamma_2} \fCenter
      \maeh{\Delta_1}{\Delta_2}$
    \DisplayProof
    \]
    இடது முன்கோளின் இடையீட்டி $!C_1$ க்காக இரண்டும், இரண்டாம்
    முன்கோளின் இடையீட்டி $!C_2$ க்காக இரண்டும் எனத் தொகுத்தறிதல்
    கருதுகோள் நான்கு !!{provable} தொடரணிகளைத் தருகிறது:
    \begin{align*}
    \Gamma_1 & \Sequent \Delta_1, !A, !C_1 &&\text{(1)} &
    !C_1, \Gamma_2 & \Sequent \Delta_2 && \text{(2)}\\
    !B, \Gamma_1 & \Sequent \Delta_1, !C_2 &&\text{(3)} &
    !C_2, \Gamma_2 & \Sequent \Delta_2 && \text{(4)}
    \end{align*}
    (1) இன் வலப்புறத்தில் $!C_2$ ஐயும் (3) இன் வலப்புறத்தில்
    $!C_1$ ஐயும் வலுவிழக்கச் சேர்த்தால், அவை~\LeftR{\lif}
    விதியின் முன்கோள்களாக அமையும்:
    \[
    \AxiomC{}
    \Deduce$\Gamma_1 \fCenter \Delta_1, !A, !C_1, !C_2$
    \AxiomC{}
    \Deduce$!B, \Gamma_1 \fCenter \Delta_1, !C_1, !C_2$
    \RightLabel{\LeftR{\lif}}
    \BinaryInf$!A \lif !B, \Gamma_1 \fCenter \Delta_1, !C_1, !C_2$
    \DisplayProof
    \]
    $\RightR{\lor}$ ஐப் பயன்படுத்தி
    \[
    !A \lif !B, \Gamma_1 \fCenter \Delta_1, !C_1 \lor !C_2
    \]
    பெறுகிறோம். ஆகவே இந்நேர்வின் இடையீட்டியாக
    $!C_1 \lor !C_2$ இருக்கலாம். எஞ்சிய (2), (4) தொடரணிகளிலிருந்து
    தேவையான மற்ற தொடரணியையும் !!{prove} செய்யலாம்:
    \[
    \AxiomC{}
    \Deduce$!C_1, \Gamma_2 \fCenter \Delta_2$
    \AxiomC{}
    \Deduce$!C_2, \Gamma_2 \fCenter \Delta_2$
    \RightLabel{\LeftR{\lor}}
    \BinaryInf$!C_1 \lor !C_2, \Gamma_2 \fCenter \Delta_2$
    \DisplayProof
    \]
    $!C_1 \lor !C_2$ சரியான மொழியில் உள்ளது என்பதை
    இன்னும் சரிபார்க்க வேண்டும். தொகுத்தறிதல் கருதுகோளால்
    \begin{align*}
    \Lang L(!C_1) & \subseteq
    \Lang L(\Gamma_1, \Delta_1, !A) \cap \Lang L(\Gamma_2, \Delta_2) &\text{மற்றும்}\\
    \Lang L(!C_2) & \subseteq
    \Lang L(!B, \Gamma_1, \Delta_1) \cap \Lang L(\Gamma_2, \Delta_2)
    \intertext{மேலும்}
    \Lang L(!A) & \subseteq \Lang L(!A \lif !B) &\text{மற்றும்}\\
    \Lang L(!B) & \subseteq \Lang L(!A \lif !B)
    \intertext{என்பதால் இரண்டிற்கும்}
    \Lang L(!C_1) & \subseteq
    \Lang L(\Gamma_1, \Delta_1, !A \lif !B) \cap \Lang L(\Gamma_2, \Delta_2) &\text{மற்றும்}\\
    \Lang L(!C_2) & \subseteq
    \Lang L(\Gamma_1, \Delta_1, !A \lif !B) \cap \Lang L(\Gamma_2, \Delta_2)
    \intertext{எனக் கிடைக்கும். இதன் விளைவாக, $\Lang L (!C_1 \lor !C_2) = {}$}
    \Lang L(!C_1) \cup \Lang L(!C_2) & \subseteq
    \Lang L(\Gamma_1, \Delta_1, !A \lif !B) \cap \Lang L(\Gamma_2, \Delta_2),
    \end{align*}
    என்று கிடைக்கிறது; இதுவே தேவை.

% SEGMENT OLP-0658-S07
    $!A \lif !B$ இரண்டாம் பகுதியில் இருந்தால் அமைப்பு மாறும்;
    இருப்பினும் இதே முறையில் கையாளலாம். தொகுத்தறிதல் கருதுகோள்
    மீண்டும் நான்கு !!{provable} தொடரணிகளைத் தருகிறது:
    \begin{align*}
    \Gamma_1 & \Sequent \Delta_1, !C_1 &&\text{(1)} &
    !C_1, \Gamma_2 & \Sequent \Delta_2, !A && \text{(2)}\\
    \Gamma_1 & \Sequent \Delta_1, !C_2 &&\text{(3)} &
    !C_2, !B, \Gamma_2 & \Sequent \Delta_2 && \text{(4)}
    \end{align*}
    (2), (4) இல் தேவையான !!{formula}s ஐ வலுவிழக்கச் சேர்த்தால்,
    அவை~\LeftR{\lif} விதியின் முன்கோள்களாக அமைகின்றன:
    \[
    \AxiomC{}
    \Deduce$!C_1, !C_2, \Gamma_2 \fCenter \Delta_2, !A$
    \AxiomC{}
    \Deduce$!B, !C_1, !C_2, \Gamma_2 \fCenter \Delta_2$
    \RightLabel{\LeftR{\lif}}
    \BinaryInf$!A \lif !B, !C_1, !C_2, \Gamma_2 \fCenter \Delta_2$
    \DisplayProof
    \]
    இங்கும் $!C_1$, $!C_2$ இலிருந்து கட்டப்படும் ஒற்றை
    !!{formula} ஒன்றை முன்பக்கத்தில் பெற வேண்டும்; ஏனெனில் இது
    இரண்டாம் பகுதி. $\LeftR{\land}$ ஐப் பயன்படுத்தினால் போதும்:
    $!C \ident (!C_1 \land !C_2)$ ஐ இடையீட்டியாகக் கொள்கிறோம்.
    எஞ்சிய (1), (3) மூலம் முதல் பகுதிக்கான தொடர்புடைய தொடரணியை
    !!{prove} செய்யலாம்:
    \[
    \AxiomC{}
    \Deduce$\Gamma_1 \fCenter \Delta_1, !C_1$
    \AxiomC{}
    \Deduce$\Gamma_1 \fCenter \Delta_1, !C_2$
    \RightLabel{\RightR{\land}}
    \BinaryInf$\Gamma_1 \fCenter \Delta_1, !C_1 \land !C_2$
    \DisplayProof
    \]
    முன்போலவே $\Lang L(!C_1 \land !C_2) \subseteq \Lang
    L(\Gamma_1, \Delta_1) \cap \Lang L(\Gamma_2, \Delta_2,
    !A \lif !B)$.

% SEGMENT OLP-0658-S08
    \item $\pi$, முதன்மை !!{formula}~$\lforall[x][!A(x)]$ உடைய
    $\RightR{\lforall}$ விதியில் முடிகிறது:
    \[
    \AxiomC{}
    \RightLabel{$\pi_1$}
    \Deduce$\maeh{\Gamma_1}{\Gamma_2} \fCenter
    \maeh{\Delta_1, !A(c)}{\Delta_2}$
    \RightLabel{\RightR{\lforall}}
    \UnaryInf$\maeh{\Gamma_1}{\Gamma_2} \fCenter
      \maeh{\Delta_1, \lforall[x][!A(x)]}{\Delta_2}$
    \DisplayProof
    \]
    அல்லது
    \[
    \AxiomC{}
    \RightLabel{$\pi_1$}
    \Deduce$\maeh{\Gamma_1}{\Gamma_2} \fCenter
      \maeh{\Delta_1}{\Delta_2, !A(c)}$
    \RightLabel{\RightR{\lforall}}
    \UnaryInf$\maeh{\Gamma_1}{\Gamma_2} \fCenter
      \maeh{\Delta_1}{\Delta_2, \lforall[x][!A(x)]}$
    \DisplayProof
    \]
    முதல் நேர்வில் தொகுத்தறிதல் கருதுகோள்
    $\Gamma_1 \Sequent \Delta_1, !A(c), !C_1$ மற்றும்
    $!C_1, \Gamma_2 \Sequent \Delta_2$ ஆகியவற்றுக்கு
    !!{proof}s தருகிறது. முதலாவதற்கு $\RightR{\lforall}$
    பயன்படுத்தி $\Gamma_1 \Sequent \Delta_1,
    \lforall[x][!A(x)], !C_1$ ஐப் பெறலாம். ($\pi$ இன்
    இறுதி \RightR{\lforall} விதியில் தனிமாறி நிபந்தனை
    ஏற்கெனவே நிறைவேறுவதால் இங்கும் அது நிறைவேறுகிறது.)
    ஆகவே $\Lang L(!C_1) \subseteq \Lang L(\Gamma_1,
    \Delta_1, \lforall[x][!A(x)]) \cap \Lang L(\Gamma_2,
    \Delta_2)$ என்ற நிபந்தனை நிறைவேறினால் $!C_1$
    இடையீட்டியாகும். தொகுத்தறிதல் கருதுகோள்
    $\Lang L(!C_1) \subseteq \Lang L(\Gamma_1, \Delta_1,
    !A(c)) \cap \Lang L(\Gamma_2, \Delta_2)$ ஐத் தருகிறது.
    எனவே $c$ பற்றி கவனிக்க வேண்டும்: $c$, $!C_1$ இல்
    வருமா? வராது. $c$ ஆனது $!C_1$ இல் வந்தால்
    $c \in \Lang L(\Gamma_1, \Delta_1, !A(c))$
    என்பதோடு $c \in \Lang L(\Gamma_2, \Delta_2)$ என்பதும்
    வேண்டும். அப்போது $\pi$ இன் $\RightR{\lforall}$ முடிவில்
    $c$ வரும்; இது தனிமாறி நிபந்தனைக்கு முரண்.

    மற்ற நேர்வும் இதேபோல் அமைகிறது.

% SEGMENT OLP-0658-S09
    \item $\pi$, முதன்மை !!{formula}~$\lexists[x][!A(x)]$
    உடைய $\RightR{\lexists}$ விதியில் முடிகிறது. இங்கும் இரு
    நேர்வுகள் உள்ளன. $\lexists[x][!A(x)]$ முதல் பகுதியில் உள்ள
    நேர்வை ஆராய்ந்து, மற்றதைப் பயிற்சியாக விடுகிறோம்.

    முதல் நேர்வில் $\pi$ பின்வருமாறு முடிகிறது:
    \[
    \AxiomC{}
    \RightLabel{$\pi_1$}
    \Deduce$\maeh{\Gamma_1}{\Gamma_2} \fCenter
      \maeh{\Delta_1, \lexists[x][!A(x)], !A(t)}{\Delta_2}$
    \RightLabel{\RightR{\lexists}}
    \UnaryInf$\maeh{\Gamma_1}{\Gamma_2} \fCenter
      \maeh{\Delta_1, \lexists[x][!A(x)]}{\Delta_2}$
    \DisplayProof
    \]
    தொகுத்தறிதல் கருதுகோளால்
    $\Gamma_1 \Sequent \Delta_1, \lexists[x][!A(x)], !A(t), !C_1$
    மற்றும் $!C_1, \Gamma_2 \Sequent \Delta_2$ ஆகியவற்றுக்கு
    !!{proof}s உள்ளன. முதலாவதற்கு $\RightR{\lexists}$ பயன்படுத்தி
    $\Gamma_1 \Sequent \Delta_1, \lexists[x][!A(x)], !C_1$
    பெறலாம். எனவே
    $\Lang L(!C_1) \subseteq \Lang L(\Gamma_1, \Delta_1,
    \lexists[x][!A(x)]) \cap \Lang L(\Gamma_2, \Delta_2)$
    என்ற நிபந்தனை நிறைவேறினால் மீண்டும் $!C_1$ இடையீட்டியாகும்.
    தொகுத்தறிதல் கருதுகோள் இங்கு
    $\Lang L(!C_1) \subseteq \Lang L(\Gamma_1, \Delta_1,
    \lexists[x][!A(x)], !A(t)) \cap \Lang L(\Gamma_2,
    \Delta_2)$ ஐத் தருகிறது.

    !!{function}s இல்லை என்ற கருதுகோளை நினைவுகொள்க. மூல
    வாதத்தின் அடுத்த கட்டம் $t$ ஐ !!a{constant} ஆகக் கொண்டு,
    $t \ident b$ என்கிறது. $b \notin \Lang L(!C_1)$ என்றாலோ,
    $b \in \Lang L(\Gamma_1, \Delta_1,
    \lexists[x][!A(x)]) \cap \Lang L(\Gamma_2, \Delta_2)$
    என்றாலோ சிக்கல் இல்லை: $!C_1$ இடையீட்டியாகும். ஆனால்
    $b \in \Lang L(!C_1)$ என்றும்
    $b \notin \Lang L(\Gamma_1, \Delta_1,
    \lexists[x][!A(x)]) \cap \Lang L(\Gamma_2, \Delta_2)$
    என்றும் இருந்தால் என்ன? $!C_1$ இல் வரும் எல்லா $b$ களுக்கும்
    பதிலாக !!a{variable}~$y$ இட்டால்,
    $!C_1 \ident \Subst{!C_1'}{b}{y}$ என்றும் $!C_1'$ இல்
    $b$ வராது என்றும் எழுதலாம். மேலும்,
    $b \notin \Lang L(\Gamma_1, \Delta_1,
    \lexists[x][!A(x)])$ அல்லது
    $b \notin \Lang L(\Gamma_2, \Delta_2)$.
    அதற்கேற்ப $!C \ident \lforall[y][!C_1'(y)]$ அல்லது
    $!C \ident \lexists[y][!C_1'(y)]$ எனக் கொள்ளலாம்.
    முதல் நேர்வில்:
    \[
    \AxiomC{}
    \Deduce$\Gamma_1 \fCenter
      \Delta_1, \lexists[x][!A(x)], !A(b), !C_1'(b)$
    \RightLabel{\RightR{\lexists}}
    \UnaryInf$\Gamma_1 \fCenter \Delta_1, \lexists[x][!A(x)], !C_1'(b)$
    \RightLabel{\RightR{\lforall}}
    \UnaryInf$\Gamma_1 \fCenter
      \Delta_1, \lexists[x][!A(x)], \lforall[y][!C_1'(y)]$
    \DisplayProof
    \]
    மேலும்
    \[
    \AxiomC{}
    \Deduce$!C_1'(b), \lforall[y][!C_1'(y)], \Gamma_2 \fCenter \Delta_2$
    \RightLabel{\LeftR{\lforall}}
    \UnaryInf$\lforall[y][!C_1'(y)], \Gamma_2 \fCenter \Delta_2$
    \DisplayProof
    \]
    மேலுள்ள தொடரணிகளின் !!{proof}s தொகுத்தறிதல் கருதுகோளால்
    கிடைக்கின்றன; இரண்டாவதில் முன்பக்கத்தில்
    $\lforall[y][!C_1'(y)]$ ஐ வலுவிழக்கச் சேர்க்கும் விதியின்
    ஏற்கத்தக்கதன்மையையும் பயன்படுத்துகிறோம். முதல் !!{proof}
    இன் \RightR{\lforall} இல் தனிமாறி நிபந்தனை நிறைவேறுகிறது:
    $b \notin \Lang L(\Gamma_1, \Delta_1,
    \lexists[x][!A(x)])$; மேலும் $!C_1'$ இல் $b$ வராதபடி
    வரையறுத்துள்ளோம்.

    இரண்டாம் நேர்வில், தொகுத்தறிதல் கருதுகோளையும்
    $\lexists[y][!C_1'(y)]$ ஐ வலுவிழக்கச் சேர்க்கும்
    ஏற்கத்தக்கதன்மையையும் கொண்டு:
    \[
    \AxiomC{}
    \Deduce$\Gamma_1 \fCenter
      \Delta_1, \lexists[x][!A(x)], !A(b), \lexists[y][!C_1'(y)], !C_1'(b)$
    \RightLabel{\RightR{\lexists}}
    \UnaryInf$\Gamma_1 \fCenter
      \Delta_1, \lexists[x][!A(x)], \lexists[y][!C_1'(y)], !C_1'(b)$
    \RightLabel{\RightR{\lexists}}
    \UnaryInf$\Gamma_1 \fCenter
      \Delta_1, \lexists[x][!A(x)], \lexists[y][!C_1'(y)]$
    \DisplayProof
    \]
    மேலும்
    \[
    \AxiomC{}
    \Deduce$!C_1'(b), \Gamma_2 \fCenter \Delta_2$
    \RightLabel{\LeftR{\lexists}}
    \UnaryInf$\lexists[y][!C_1'(y)], \Gamma_2 \fCenter \Delta_2$
    \DisplayProof
    \]
    இரண்டாம் !!{proof} இன் \LeftR{\lexists} இல்
    தனிமாறி நிபந்தனை நிறைவேறுகிறது: $b \notin \Lang
    L(\Gamma_2, \Delta_2)$; மேலும் $!C_1'$ இல் $b$ வராது.
  \end{enumerate}
  மற்ற நேர்வுகளும் இதேபோல அமைவதால் அவை பயிற்சிகளாக
  விடப்படுகின்றன.
\end{proof}

பதிப்புக் குறிப்பு: சார்புக்குறியீடுகள் இல்லாவிட்டாலும் மாறிகள்
பதங்களாக இருக்கலாம். மேலுள்ள மூல வாதம் மாறிலி சாட்சி என்ற
மாறிலி நேர்வை மட்டுமே விரிவாகக் கையாள்கிறது; மாறி நேர்வுக்கான
கூடுதல் படி உறைந்த மூலத்தில் வழங்கப்படவில்லை.

% SEGMENT OLP-0658-S10
\begin{prob}
$\lnand$ என்ற இணைப்பிக்கு பின்வரும் விதிகள் இருப்பதாகக் கொள்க:
\[
\Axiom$\Gamma \fCenter \Delta, !A$
\Axiom$\Gamma \fCenter \Delta, !B$
\RightLabel{\LeftR{\lnand}}
\BinaryInf$!A \lnand !B, \Gamma \fCenter \Delta$
\DisplayProof
\quad
\Axiom$!A, !B, \Gamma \fCenter \Delta$
\RightLabel{\RightR{\lnand}}
\UnaryInf$\Gamma \fCenter \Delta, !A \lnand !B$
\DisplayProof
\]
$\pi$ இவ்விதிகளில் ஒன்றில் முடியும் நேர்வுகளை
\olref{lem:maehara} இன் நிறுவலில் செய்து,
மாஎஹாராவின் துணைத்தேற்றம் தொடர்ந்து பொருந்துகிறது எனக் காட்டுக.
\end{prob}

மாஎஹாராவின் துணைத்தேற்றத்தை நிறுவியதால், அதைக் கொண்டு கிரெய்கின்
இடையீட்டுத் தேற்றத்தை நிறுவலாம்.

% SEGMENT OLP-0658-S11
\begin{proof}[\olref{thm:interpolation} இன் நிறுவல்]
  $!A \Sequent !B$ !!{provable} என்று கொள்க.
  $\maeh{!A}{\ } \Sequent \maeh{\ }{!B}$ க்கு
  \olref{lem:maehara} ஐப் பயன்படுத்தி இடையீட்டி~$!C$
  பெறுக. அப்போது $\Proves !A \Sequent !C$ மற்றும்
  $\Proves !C \Sequent !B$.
\end{proof}

% SEGMENT OLP-0658-S12
$n$-இடப் பயனிலை~$R$ ஐக் கொண்ட !!a{language}~$\Lang L$
இல் $\Gamma$ ஒரு கோட்பாடு, அதாவது !!{sentence}s இன் கணம் என்க.
பின்வரும் நிபந்தனை இருந்தால், இருந்தால்தான், $\Gamma$ ஆனது
$R$ ஐ \emph{மறைமுகமாக வரையறுக்கிறது} என்போம்:
\begin{align*}
\Gamma, \Gamma'
& \Entails \lforall[x_1][\dots\lforall[x_n][(R(x_1,\dots,x_n) \liff
R'(x_1, \dots, x_n))]],
\intertext{இங்கு $\Gamma'$ என்பது $\Gamma$ இல் வரும் $R$ இன்
எல்லா நிகழ்வுகளுக்கும் பதிலாக $R' \notin \Lang L$ ஐ இட்டுப்
பெறப்பட்டது. $\Gamma$ ஆனது $R$ ஐ
\emph{வெளிப்படையாக வரையறுக்கிறது} என்பதற்கு,
$!A(x_1, \dots, x_n)$ என்ற, $R$ இல்லாத வாய்பாடு ஒன்று இருந்து
பின்வருவது நிறைவேற வேண்டும்:}
\Gamma &\Entails
\lforall[x_1][\dots\lforall[x_n][(R(x_1,\dots,x_n) \liff !A(x_1,
\dots, x_n))]].
\end{align*}
$\Gamma$ ஆனது $R$ ஐ வெளிப்படையாக வரையறுத்தால் $R$ ஐ மறைமுகமாகவும்
வரையறுக்கும் என்பது தெளிவு. இதன் மறுதிசை உடனடியாகத் தெளிவாகாவிட்டாலும்
உண்மையாகும்:

% SEGMENT OLP-0658-S13
\begin{lem}[பெத் வரையறுக்கத்தக்கமைத் தேற்றம்]
  $\Gamma$ ஆனது $R$ ஐ மறைமுகமாக வரையறுத்தால்,
  அதனை வெளிப்படையாகவும் வரையறுக்கிறது.
\end{lem}

\begin{proof}
  $\Gamma$ ஆனது $R$ ஐ மறைமுகமாக வரையறுக்கிறது என்க.
  மேலே கூறிய $R'$, $\Gamma'$ ஆகியவற்றை எடுத்துக்கொண்டு
  $\Lang L' = (\Lang L \setminus \{R\}) \cup \{R'\}$
  என்க. கருதுகோளால்,
  \begin{align*}
    \Gamma, \Gamma' &
    \Sequent \lforall[x_1][\dots\lforall[x_n][(R(x_1,\dots,x_n)
    \liff R'(x_1, \dots, x_n))]]
  \intertext{$c_1$, \dots,~$c_n$ என்பவை $\Gamma$ இல் வராத
  புதிய !!{constant}s என்க. மீளாக்கத் துணைத்தேற்றத்தால்}
    \Gamma, \Gamma', R(c_1,\dots,c_n) & \Sequent  R'(c_1, \dots, c_n)
  \intertext{பெறுகிறோம். பின்வருவதற்கு மாஎஹாராவின் துணைத்தேற்றத்தைப்
  பயன்படுத்துக:}
    \maeh{\Gamma, R(c_1,\dots,c_n)}{\Gamma'} & \Sequent
      \maeh{\ }{R'(c_1, \dots, c_n)}
  \intertext{கிடைக்கும் இடையீட்டியில் புதிய மாறிலிகளை முறையே
  கட்டிலா மாறிகளால் மாற்றி $!A \ident !A(x_1, \dots, x_n)$
  பெறுகிறோம். அப்போது பின்வரும் தொடரணிகளுக்கு !!{proof}s உள்ளன:}
    \Gamma, R(c_1,\dots,c_n) & \Sequent !A(c_1, \dots, c_n) \\
    !A(c_1, \dots, c_n), \Gamma' & \Sequent R'(c_1, \dots, c_n)
  \intertext{இடையீட்டியின் சொற்குறிகள் இருபுறத்துக்கும் பொதுவானவை.
  புதிய மாறிலிகளை மாறிகளால் மாற்றிய பிறகு,
  $\Lang L(!A) \subseteq \Lang L(\Gamma) \setminus \{R\}$;
  எனவே $!A$ இல் $R$ வராது. இரண்டாம் தொடரணியின் !!{proof}
  முழுவதிலும் $R'$ க்குப் பதிலாக $R$ ஐ இட்டால் பின்வருவதற்கான
  !!a{proof} கிடைக்கும்:}
    !A(c_1, \dots, c_n), \Gamma & \Sequent R(c_1, \dots, c_n)
  \end{align*}
  இரு திசைகளின் விளைவுகளையும் சேர்த்து, \RightR{\lif} மற்றும்
  \RightR{\lforall} ஐப் பயன்படுத்தினால் பின்வருவதற்கான
  !!a{proof} கிடைக்கும்:
  \[
    \Gamma \Sequent \lforall[x_1][\dots\lforall[x_n][(R(x_1,\dots,x_n) \liff !A(x_1, \dots, x_n))]].
  \]
  ஆகவே $!A(x_1, \dots, x_n)$, $\Gamma$ இல் $!R$ ஐ
  வெளிப்படையாக வரையறுக்கிறது.
\end{proof}

% SEGMENT OLP-0658-S14
இடையீட்டுக்குச் சமானமானதும், அதைப் போல மாஎஹாராவின்
துணைத்தேற்றத்தால் நிறுவக்கூடியதுமான மூன்றாம் முடிவு இதுவாகும்:
முரணின்மையுடைய $\Gamma_1$, $\Gamma_2$ என்ற இரு கோட்பாடுகளும்,
அவற்றின் பொதுமொழியான
$\Lang L(\Gamma) = \Lang L(\Gamma_1) \cap \Lang L(\Gamma_2)$
இல் நிறைவான $\Gamma$ என்ற கோட்பாட்டை உட்கொண்டால்,
$\Gamma_1 \cup \Gamma_2$ முரணின்மையுடையது.

நிறைவான கோட்பாட்டின் மொழியில் உள்ள ஒவ்வொரு
!!{sentence}~$!A$ க்கும் $\Gamma \Proves !A$ அல்லது
$\Gamma \Proves \lnot !A$ ஆகும். $\Gamma_1$ உம்
$\Gamma_2$ உம் $\Gamma$ இன் முரணின்மையுடைய நீட்சிகள் என்பதால்
$\Gamma$ உம் முரணின்மையுடையது. ஆகவே
$\Lang L(\Gamma_1) \cap \Lang L(\Gamma_2)$ இன்
!!a{sentence}~$!A$ க்கு ஒரே நேரத்தில்
$\Gamma_1 \Entails !A$ மற்றும்
$\Gamma_2 \Entails \lnot !A$ ஆக முடியாது. அவ்வாறு
ஒரு $!A$ இருப்பதாகக் கொள்க. $\Gamma_1 \Entails !A$
எனில் $\Gamma \Entails !A$ ஆகவேண்டும்: இல்லையேல்
$\Gamma$ இன் நிறைவுமையால் $\Gamma \Entails \lnot !A$;
$\Gamma \subseteq \Gamma_1$ என்பதால்
$\Gamma_1 \Entails \lnot !A$ ஆகி $\Gamma_1$
முரண்படும். எனவே $\Gamma \Entails !A$; மேலும்
$\Gamma \subseteq \Gamma_2$ என்பதால்
$\Gamma_2 \Entails !A$. இப்போது
$\Gamma_2 \Entails/ \lnot !A$; இல்லையெனில்
$\Gamma_2$ முரண்படும்.

பொது !!{language}~$\Lang L(\Gamma_1) \cap
\Lang L(\Gamma_2)$ இல் $\Gamma_1 \Entails !A$ மற்றும்
$\Gamma_2 \Entails \lnot !A$ ஆகிய இரண்டும் நிறைவேறும்
!!a{sentence}~$!A$ இல்லாமை மட்டுமே நிறுவலுக்குத் தேவை.
மாஎஹாராவின் துணைத்தேற்றத்தைக் கொண்டு நிறுவல்-கோட்பாட்டு முறையில்
அதை இப்போது நிறுவுகிறோம்.

% SEGMENT OLP-0658-S15
\begin{thm}[ராபின்சனின் கூட்டு முரணின்மைத் தேற்றம்]
$\Gamma_1$, $\Gamma_2$ முரணின்மையுடைய கோட்பாடுகள் என்க.
!!{language}~$\Lang L(\Gamma_1) \cap \Lang L(\Gamma_2)$ இல்
உள்ள எந்த $!A$ க்கும் $\Gamma_1 \Entails !A$ மற்றும்
$\Gamma_2 \Entails \lnot !A$ ஆகிய இரண்டும் நிறைவேறவில்லை
எனில், $\Gamma_1 \cup \Gamma_2$ முரணின்மையுடையது.
\end{thm}

\begin{proof}
  மாறாக $\Gamma_1 \cup \Gamma_2$ முரண்படுகிறது என்க.
  அப்போது $\Gamma_1, \Gamma_2 \Sequent \ $ என்ற தொடரணிக்கு
  !!a{proof} உள்ளது. $\maeh{\Gamma_1}{\Gamma_2}
  \Sequent \maeh{\ }{\ }$ க்கு மாஎஹாராவின் துணைத்தேற்றத்தைப்
  பயன்படுத்தினால், பொது !!{language}~$\Lang L(\Gamma_1) \cap
  \Lang L(\Gamma_2)$ இல் !!a{sentence}~$!A$ ஒன்று கிடைக்கும்;
  அதற்கு $\Proves \Gamma_1 \Sequent !A$ மற்றும்
  $\Proves !A, \Gamma_2 \Sequent \quad$.
  இரண்டாவதிலிருந்து $\Proves \Gamma_2 \Sequent \lnot !A$
  பெறுகிறோம். அத்தகைய !!{sentence} இல்லை என்று கருதினோம்.
\end{proof}

மேலே $\Gamma$, $\Gamma_1$, $\Gamma_2$ முடிவுறு கோட்பாடுகள்
என்று மறைமுகமாகக் கொண்டோம். கச்சிதத்தன்மையால் முடிவிலாக்
கோட்பாடுகளுக்கும் முடிவுகள் பொருந்தும்.

\end{document}
